\documentclass{article}
\usepackage{graphicx} % Required for inserting images

\title{Ogres\_villagers\_and\_the\_art\_of\_writing}
\author{Rajat Ranjan Meher\\[0.2cm]
\textbf{Instructor:-} Bhaskaran Raman   
}
\date{Even semseter: 2026}

\begin{document}

\maketitle

\section{Summary}

\subsection{Introduction}
This paper is about increasing effective communication of ideas using an analogy of age old "Ogres and villagers" story. It has mainly discussed the shortcoming authors make while writing a paper and how avoid it and gives an approach to move toward better communication.
\subsection{Content}
In this paper author has discussed four steps.\\

\textbf{First:-} Thoroughly describe the problem statements and all the player in your research.As in the story who are villagers and ogres and don't leave it to reader's imagination.\\

\textbf{Second:-} Explore all previous solution attempted and analyze in a convincing way to show why they work or did not work and what may be it's short comings.Describe in sufficient detail that other too can perform the solution.   \\

\textbf{Third:-} Walk the reader through your solution in details. In context of the story "What makes it ogre proof? What materials did you use? How much did itcost? How long did it take to build?"\\

 \textbf{Conclusion:-}Show them why it works or up to what extend it works. As in the story did it stop the ogre or at least delay their movements. Draw a statistically sound conclusion. \\

\subsection{Consequences:-}
Writing a paper in an poor way can lead to frustration of the reviewer ,who then can retaliate with low scores and summary rejection. The most common complaint reviewer gave is lack of sufficient explanation of solution. Directly jumping into solution can lead to an idea that you are simply reinventing the wheel. 

\end{document}

